\section{Data}
\label{sec:data}

Our system draws on six data sources: the official FIFA Fantasy API
(the primary inference source), a public archive of international
match results \citep{martj42}, club-level match and odds data
\citep{footballdata}, community dumps of English Premier League
fantasy scoring \citep{vaastav}, a static copy of the FIFA Men's
World Ranking as a fallback, and a community-maintained World Cup
2026 statistics dataset \citep{mominul2026dataset} providing real
match expected goals, lineups with minutes, and referee card
histories. This section describes each source, the schemas we
consume, and the harmonisation work required to join them.

\subsection{FIFA Fantasy API}
\label{sec:data-fifa-api}

The primary inference source consists of three JSON endpoints under
\texttt{play.fifa.com/json/fantasy/}:

\begin{itemize}
\item \texttt{players.json}: one record per player, with \texttt{id},
  position, price (in millions), \texttt{percentSelected} (ownership,
  0--100), \texttt{status} (playing, transferred, and similar),
  \texttt{oneToWatch}, and a \texttt{stats} block containing
  \texttt{totalPoints}, \texttt{lastRoundPoints}, \texttt{form},
  \texttt{roundPoints}, and \texttt{nextFixtureFromActiveRound}.
\item \texttt{squads.json}: one record per national team, with
  \texttt{id}, \texttt{name}, \texttt{abbr}, \texttt{group}, and
  \texttt{isEliminated}.
\item \texttt{rounds.json}: one record per round, with status, start
  and end dates, and a \texttt{tournaments} array of fixtures (home
  and away squad ids, kickoff time, and scores once played).
\end{itemize}

\paragraph{Schema drift at MD1 kickoff.}
Pre-tournament, \texttt{stats.roundPoints} was a list of integers. On
the day Matchday~1 started, the field switched to a dictionary keyed
by round id (for example \texttt{\{"1": 7\}}). Our Pydantic schema
initially modelled the field as a list, which broke the parser at the
worst possible moment. The fix accepts both shapes via a
\texttt{list[int] | dict[str, int]} union type and a normaliser that
emits a list indexed by round position, so downstream code is
insulated from the representation.

\paragraph{Two-stage validation.}
A first layer of raw Pydantic models mirrors the API verbatim
(camelCase field names, all-optional wherever the source is
unreliable). A second normalisation pass produces the downstream
contract (\texttt{Squad}, \texttt{Player}, \texttt{Fixture}) consumed
by every later module. This is the standard anti-corruption-layer
pattern: the API can drift without breaking the rest of the codebase,
provided only the normalisation pass is updated.

\subsection{International match results and live Elo}
\label{sec:data-elo}

We consume a public CSV of every international match since 1872
\citep{martj42}, with columns \texttt{date}, \texttt{home\_team},
\texttt{away\_team}, \texttt{home\_score}, \texttt{away\_score},
\texttt{tournament}, \texttt{city}, \texttt{country}, and
\texttt{neutral}. The companion goalscorer and shootout files exist
but we do not consume them.

From this history we compute an Elo rating \citep{elo1978rating} per
country by rolling forward through the date-sorted match list. The
K-factor varies by tournament weight: 60 for World Cup matches, 25
for qualifiers, and 10 for friendlies. A goal-margin adjustment
scales the K-factor by $\sqrt{(m+1)/2}$ for a margin of $m$ goals,
and home advantage adds 60 to the home side's expected score except
in neutral fixtures; both choices follow common practice in
Elo-based football forecasting \citep{hvattum2010using}. The output
table holds, per country, the current Elo, match count, form over the
last ten matches, goals for and against over the trailing 24 months,
and the date of the most recent match.

\subsection{Club-level enrichment}
\label{sec:data-clubs}

We download per-league, per-season CSVs from football-data.co.uk
\citep{footballdata} for the top five European leagues (Premier
League, La Liga, Bundesliga, Serie A, Ligue 1), covering the most
recent three full seasons plus earlier seasons used only as Elo seed
data. The data serves three purposes:

\begin{itemize}
\item a per-club Elo rating, computed analogously to the
  international Elo above;
\item a normalised match table carrying bookmaker closing odds
  (Pinnacle, with Bet365 as fallback) where available;
\item a per-club, per-date Elo timeline used to attach
  team-Elo-at-match-time to the Premier League training rows without
  lookahead bias.
\end{itemize}

\subsection{Premier League fantasy dumps as training data}
\label{sec:data-fpl}

Labelled training data comes from a public mirror of every English
Premier League fantasy round \citep{vaastav}. We ingest three full
seasons (2022--23, 2023--24, 2024--25) into per-player, per-gameweek
Parquet tables. After dropping non-appearances (zero-minute rows), we
retain \EplTrainingRows{} labelled rows across
the three seasons.

\subsection{Static FIFA World Ranking fallback}
\label{sec:data-rankings}

We also maintain a static copy of the FIFA Men's World Ranking,
hand-transcribed from FIFA's own publication. It serves as a fallback
for the country-strength feature when the live Elo of
Section~\ref{sec:data-elo} is unavailable, mostly for countries that
have not played any international match within the cached window. The
heuristic and Poisson backends prefer Elo over the static rank when
both are present, fall back to the rank when Elo is missing, and fall
back to a price-only signal when neither is present.

\subsection{Country-name harmonisation}
\label{sec:data-names}

A persistent integration cost is that each source spells country
names differently: United States versus USA versus United States of
America, or Korea Republic versus South Korea. A hand-maintained
mapping table holds canonical lookups in two directions: one mapping
used when joining the international Elo into the per-player feature
table, and one mapping club names between the fantasy dumps and the
club-level match data. Three near-misses had to be fixed during
integration: Bosnia and Herzegovina (both sources agree, but our
initial mapping incorrectly rewrote it); Czechia (FIFA) versus Czech
Republic in the match archive; and T\"urkiye (FIFA, the country's
official name since 2022) versus Turkey. Each was caught by an
after-join NaN scan over the country-Elo column of the per-player,
per-round feature table.
